% Three-page Korean resume. Personal identifiers omitted except date of birth.
% Build from the repository root: xelatex resume_brief_3page_korean.tex
\documentclass[a4paper,11pt]{article}
\usepackage{fontspec}
\usepackage{kotex}
\setmainfont{NanumGothic}[Path=fonts/,Extension=.ttf,UprightFont=*,BoldFont=*Bold]
\usepackage[margin=19mm,top=18mm,bottom=19mm]{geometry}
\usepackage{xcolor}
\usepackage{enumitem}
\usepackage{titlesec}
\usepackage{tabularx}
\usepackage{fancyhdr}
\usepackage[unicode,pdfauthor={},pdftitle={경력 요약 이력서}]{hyperref}
\definecolor{accent}{HTML}{193D50}
\definecolor{muted}{HTML}{536570}
\pagestyle{fancy}
\fancyhf{}
\fancyfoot[L]{\footnotesize\color{muted}Data Engineering · Selected Projects}
\fancyfoot[R]{\footnotesize\color{muted}\thepage\ / 3}
\renewcommand{\headrulewidth}{0pt}
\setlength{\parindent}{0pt}
\setlength{\parskip}{5pt}
\setlength{\tabcolsep}{0pt}
\linespread{1.13}
\raggedright
\titleformat{\section}{\large\bfseries\color{accent}}{}{0pt}{}[\vspace{2pt}\titlerule]
\titlespacing*{\section}{0pt}{13pt}{7pt}
\setlist[itemize]{leftmargin=14pt,itemsep=4pt,topsep=3pt,parsep=0pt}
\newcommand{\role}[3]{\textbf{#1}\hfill{\small #2}\\{\small\color{muted}#3}\par}
\newcommand{\project}[2]{\vspace{5pt}\textbf{\color{accent}#1}\hfill{\small\color{muted}#2}\par}
\newcommand{\birthdate}{1992.01.02}
\begin{document}
{\fontsize{25}{31}\selectfont\bfseries\color{accent}데이터 엔지니어 프로젝트 이력서}\par
{\small 생년월일: \birthdate}\par
새로운 도메인을 학습하고 연구자 인터뷰로 문제를 정의해, 사내 MLOps를 처음부터 설계·구축했습니다. 요구사항 발굴과 기술 선정부터 하드웨어 구성·데이터 파이프라인 운영까지 직접 주도합니다.

\section{HUINNO | 연구 데이터 파이프라인 및 레이크하우스}
\role{Data Engineer}{2025.12 - 현재}{담당 범위: ETL·레이크하우스 및 온프레미스 하드웨어·Kubernetes 구축·운영 주도}

\project{01. 요구사항 발굴부터 사내 MLOps 신규 설계·구축}{2025.12 - 현재}
\textbf{과제} · 의료 데이터 도메인과 온프레미스 MLOps 구축을 새롭게 익히며, 담당자별 스크립트와 파일 전달에 의존하던 연구 업무를 지원할 공용 플랫폼을 처음부터 만들어야 했습니다.
\begin{itemize}
\item \textbf{문제 정의·요구사항 발굴:} 연구자 인터뷰와 데이터·코드·실행 흐름 분석을 직접 수행했습니다. 데이터 준비의 특정 담당자 의존, 분산된 실험 기록, 버전·생성 경로 추적 부재를 파악해 자동화·공용 저장·재현성을 설계 요구사항으로 정리했습니다.
\item \textbf{요구사항에 따른 기술 선정:} 작업 실행·재시도는 Airflow, SQL 변환은 dbt, 테이블 버전 관리는 Iceberg, 출처 추적은 DataHub로 역할을 나눴습니다. 사내 장비의 공용 실행 환경으로 Kubernetes를, S3 호환성과 NVMe/HDD 계층화를 위해 SeaweedFS를 선택했습니다.
\item \textbf{설계에서 구현·운영까지:} 수집 → Bronze/Silver/Gold 변환 → 계보 → SQL 조회를 연결하고, 공용 ETL 프레임워크와 제품·MIMIC-IV·유지보수 파이프라인을 구현했습니다. PostgreSQL + pg\_duckdb로 연구자의 분석 조회 환경도 구성했습니다.
\item \textbf{연구자와 역할·도입 방향 조율:} 도메인 변환은 연구자가 확장하고 실행·계보·검증·재시도는 플랫폼이 맡도록 경계를 정의했습니다. 구축 후에도 연구자 인터뷰로 사용상의 병목을 정리하고 제품별 도입 로드맵을 구체화했습니다.
\end{itemize}
\textbf{결과} · 도메인 학습·문제 정의·기술 선정에서 구현·운영까지 주도해 사내 공용 데이터 플랫폼을 구축했습니다. 연구자가 직접 데이터를 조회·활용하고 재현 기능을 확장할 기반을 마련했습니다.

\project{02. 온프레미스 베어메탈 Kubernetes 구축·운영}{}
\textbf{문제} · 사내 보유 장비로 데이터 플랫폼을 운영해야 했으며, 장비별 사양 차이와 메모리·스토리지·네트워크 병목을 함께 해결해야 했습니다.
\begin{itemize}
\item \textbf{온프레미스 구축 주도:} 플랫폼 2인 팀에서 사내 물리 장비를 활용한 베어메탈 Kubernetes 클러스터 구성을 단독 담당했습니다. 12개 Pulumi 스택으로 인프라를 코드화하고 10노드 규모로 운영했습니다(2026.09 기준).
\item \textbf{하드웨어 재구성·증설 주도:} 유휴 PC·워크스테이션에서 RAM·NVMe 등 부품을 직접 회수·재장착해 기존 노드의 자원을 확충했습니다. 장비 구성부터 메모리·디스크 증설, 네트워크 장비 선정·교체까지 직접 수행했습니다.
\item \textbf{시스템 자원 진단:} NUMA·듀얼채널 구성과 메모리 부족 시 SSH·Kubernetes 영향을 분석했습니다. Swap 설정과 PostgreSQL 데이터 디스크 분리로 운영 설정을 조정했습니다.
\item \textbf{저장·전송 경로 개선:} SeaweedFS 기반 Iceberg 저장소의 NVMe/HDD 계층을 구성했습니다. 케이블·스위치까지 데이터 전달 경로를 점검하고 필요한 장비 교체와 저장소 설정을 직접 수행했습니다.
\end{itemize}
\textbf{결과} · 사내 유휴 부품을 운영 자원으로 전환하고, 물리 장비부터 OS·Kubernetes·데이터 저장소까지 연결해 ETL과 분석을 지원하는 온프레미스 운영 기반을 구축했습니다.

\newpage
{\fontsize{21}{27}\selectfont\bfseries\color{accent}데이터 신뢰성 및 분산 처리}\par
\section{HUINNO | Iceberg 기반 재현 가능한 MLOps 구성}
\role{Data Engineer}{2025.12 - 현재}{담당 범위: 공용 저장소·데이터 계보·버전 관리·데이터 동결 및 복원 기반}
\project{03. 데이터 신뢰성·버전 관리·재현 기반 구축}{구축·운영 및 재현 기능 확장}
\textbf{문제} · 데이터가 노드별로 흩어져 있고 생성·저장 방법은 소수 생산자의 암묵지에 의존해, 담당자 부재 시 필요한 데이터를 요청·확보하기 어려웠습니다. 생성 코드도 버전 관리되지 않아 같은 데이터를 재현할 기반이 없었습니다. 데이터는 SSH 접속 후 파일 복사로 전달하고, 모델과 성능은 Notion에서 공유·관리해 데이터·코드·모델·성능을 함께 추적하기 어려웠습니다.
\begin{itemize}
\item \textbf{공용 저장소·카탈로그:} 노드별 데이터를 SeaweedFS의 S3 호환 저장소로 통합하고 Iceberg REST Catalog로 Silver/Gold 테이블의 데이터·메타데이터를 관리했습니다. Iceberg 스냅샷·버전과 DataHub 계보를 데이터 재현의 기반으로 구성했습니다.
\item \textbf{분석 조회 계층:} PostgreSQL + pg\_duckdb로 Iceberg SQL 조회 환경을 구성했습니다. 카탈로그의 최신 메타데이터 경로를 DuckDB에 전달해 Silver dimension을 생성하도록 했습니다.
\item \textbf{데이터 동결·복원:} manifest로 데이터 버전을 고정하고 동일 데이터를 복원하는 공용 도구를 구현했습니다. 복원·재현 정보 조회 기능과 사용 가이드를 반영하고 자동 검사를 통과했습니다.
\item \textbf{실험·모델 추적 기반:} MLflow 트래킹·모델 레지스트리와 SeaweedFS 아티팩트 저장소를 구성했습니다. 데이터 버전·코드·실행 환경을 실험 기록과 연결하는 제품별 연구 코드 연동을 진행하고 있습니다.
\item \textbf{운영 신뢰성:} 누락 원천 파일의 재수집과 무결성 확인을 수행했습니다. Airflow·DB는 서비스 단위로, 저장소는 검증 전용 버킷·접근 권한으로 격리해 운영 데이터에 대한 영향을 통제했습니다.
\end{itemize}
\textbf{결과} · 카탈로그·데이터 버전·계보·동결 및 복원을 연결해 동일한 학습 데이터를 재사용할 수 있는 MLOps 기반을 구축했습니다. 제품 연구 코드 연동과 전체 데이터·모델 재현 검증은 진행 중입니다.

\section{Eazel | 미술 시장 통합 데이터 플랫폼}
\role{Team Lead, AI Labs}{2022.06 - 2024.09}{담당 범위: 데이터 파이프라인 설계·구현·배포·모니터링·최적화 총괄}
\project{04. AWS EMR/PySpark 기반 정제·적재 파이프라인}{2023.01 - 2023.10}
\textbf{문제} · 작가·작품·경매·기관 등 여러 출처의 데이터를 통합해 미술품 분석 서비스에 제공해야 했습니다. 수집량 증가를 고려한 처리 구조와 반복 가능한 정제·적재 과정이 필요했습니다.
\begin{itemize}
\item \textbf{처리 구조 설계:} 팀 리드로 통합 파이프라인을 설계하고 AWS EMR/PySpark 정제·적재와 분산 처리를 위한 파티셔닝을 구현했습니다.
\item \textbf{처리 단계 구성:} 각 처리 단계를 Python 파일로 정의하고 EMR 클러스터 생성 시 EMR Step으로 순차 실행하도록 구성했습니다.
\item \textbf{워크플로우 연동:} Airflow DAG에서 boto3로 EMR 작업을 실행하도록 연결해 데이터 수집 이후의 정제·적재를 주기적으로 수행했습니다.
\end{itemize}
\textbf{결과} · 여러 출처의 데이터를 정제·적재하는 미술 시장 분석 서비스의 기반을 구축했습니다. 데이터 규모 증가에 대응할 수 있는 분산 처리 구조와 주기 실행 체계를 마련했습니다.



\newpage
{\fontsize{21}{27}\selectfont\bfseries\color{accent}수집 운영·데이터 품질·모델링}\par
\section{Eazel | 수집부터 분석용 데이터 제공까지}
\role{Team Lead, AI Labs}{2022.06 - 2024.09}{담당 범위: 수집 시스템 운영·중복 제거 최적화·웨어하우스 스키마 설계}
\project{05. 데이터 수집 인스턴스 중앙 관리 및 실행 자동화}{2023.01 - 2024.02}
\textbf{문제} · 50개 이상의 스크래퍼 인스턴스를 개별 관리하면서 상태 확인과 장애 대응에 운영 부담이 발생했습니다. 수집 실행부터 후속 처리까지 중앙에서 관리할 필요가 있었습니다.
\begin{itemize}
\item \textbf{중앙 제어:} Flask 기반 관리 API와 WebSocket 연결을 구현해 인스턴스에 수집 대상을 전달하고 상태·메트릭을 교환하도록 구성했습니다.
\item \textbf{실행 자동화:} 정적·동적 웹/API 수집기를 개발하고 Airflow에서 관리 API를 호출해 수집을 주기 실행했습니다. 후속 정제·적재는 EMR 작업으로 연결했습니다.
\end{itemize}
\textbf{결과} · 50개 이상의 수집 인스턴스를 한곳에서 제어하고 실행 상태와 로그를 중앙에서 확인하도록 개선했습니다. 문제가 발생한 인스턴스를 식별하고 원격 재시작할 수 있는 운영 체계를 구축했습니다.

\project{06. pgvector 기반 이미지·텍스트 중복 제거 최적화}{2023.11 - 2024.02}
\textbf{문제} · 기존 Faiss-CPU 방식은 PostgreSQL에서 임베딩을 내려받고 데이터 타입을 변환한 뒤 외부 인덱스를 생성했습니다. 반복적인 전송·변환·인덱스 생성이 처리 부담을 만들었습니다.
\begin{itemize}
\item \textbf{분석·협업:} ML 리서치 엔지니어와 구축한 초기 중복 제거 시스템의 오버헤드를 분석하고 PostgreSQL 내부에서 벡터 검색을 수행하는 구조로 개선했습니다.
\item \textbf{임베딩 활용:} Facebook SSCD 이미지 임베딩과 OpenAI 텍스트 임베딩을 활용해 유사 데이터를 탐색하도록 구성했습니다.
\item \textbf{저장·검색 변경:} 기존 NumPy 배열을 pgvector의 vector 타입으로 전환하고 IVFFlat·HNSW 인덱스를 적용했습니다.
\end{itemize}
\textbf{결과} · 외부로 데이터를 전송·변환하는 과정을 줄이고 중복 제거 구조를 단순화했습니다. 인덱스 생성 오버헤드를 50\% 줄였습니다.

\project{07. SSoT 기반 데이터 웨어하우스 스키마 재설계}{2024.02 - 2024.09}
\textbf{문제} · 단일 RDS가 원천 저장과 웨어하우스를 겸하면서 중복 데이터 관리와 비용 부담이 커졌습니다. 여러 출처의 동일 데이터를 일관되게 관리할 저장 구조가 필요했습니다.
\begin{itemize}
\item \textbf{계층 분리 설계:} S3를 데이터 레이크로 활용하는 방향으로 DataSource·Data Warehouse·Data Mart 계층을 재설계했습니다.
\item \textbf{기준 데이터 모델:} 여러 출처에서 들어온 중복 데이터를 통합 관리할 SSoT(Single Source of Truth) 테이블을 포함해 약 50개 테이블의 스키마를 설계했습니다.
\end{itemize}
\textbf{결과} · 약 50개 테이블의 스키마로 중복 데이터의 기준과 저장 계층별 역할을 정의했습니다. S3 데이터 레이크 도입과 RDS 역할 분리를 위한 구조·스키마 설계를 마련했습니다.


\section{자격·어학}
\textbf{정보처리기사} · 2017.11 취득\hfill\textbf{영어 OPIc AL}
\end{document}
